%%%%%%%%%%%%%%%%%%%%%%%%%%%%%%%%%%%%%%%%%%%%%%%%%%%%%%%%%%%%%%%%%%%%%%%%%%%%%%%%
\chapter{Introduction}\label{ch:introduction}
%%%%%%%%%%%%%%%%%%%%%%%%%%%%%%%%%%%%%%%%%%%%%%%%%%%%%%%%%%%%%%%%%%%%%%%%%%%%%%%%

Functional programming, with its foundation being the $\lambda$-calculus,
enjoys a plethora of strong, theoretical properties. Particularly with
the addition of structured data in the form of data types, many functional programming
languages allow for reasoning about complex flow of data.

However, with the surge of mobile applications, web services, and data centers,
the nature of programming is becoming increasingly parallel. Thus, our programs
need not only structure their data, but also their communication.
Consequently, process calculi emerged in the search of a model for concurrent
computation---Milner's $\pi$-cal\-cu\-lus \cite{Milner1992:PiCalculus} being a prominent example.

The need for structured concurrency gave rise to \emph{session types}
which have been introduced by Honda \cite{Honda1993:TypesForDyadicInteraction}.
Just as data types encode the structure
of our data, session types specify concurrent interaction as type-theoretic
communication protocols. Approaching session types as a static type system for $\pi$-calculi,
we would like for these protocols to be enforced at compile time with
two safety guarantees being of particular interest: type preservation and
deadlock freedom. One way for a type system to satisfy these safety guarantees
is to approach its design through the relation to formal logic.

\section{Curry-Howard: Proofs as Programs}
One of the advantages of functional programming languages is their tight correspondence
to mathematical logic. The \emph{Curry-Howard correspondence} \cite{Curry1959-CombinatoryLogic, Howard1980-FormulasAsTypes}
provides the theoretical foundation relating \emph{proofs} and \emph{programs},
\emph{propositions} and \emph{types}.

Most famously, this connection is explained
through Church's $\lambda$-calculus \cite{Church1932:LambdaCalc} and
Gentzen's natural deduction $\ND$ \cite{Gentzen1935:Untersuchungen1}:
Proofs in intuitionistic logic correspond to programs written in the $\lambda$-calculus,
and, vice versa, terms of the $\lambda$-calculus give an algorithmic specification
of an intuitionistic proof. But the correspondence goes beyond a mere bijection
between proofs and programs; it also connects the rules for evaluating programs
to the ones for simplifying proofs.

Natural deduction is not the only logic that witnesses this connection. The
Curry-Howard view has proved to be more general and has been applied to various
other logics, such as Gentzen's sequent calculus $\LK$ \cite{Gentzen1935:Untersuchungen2}---introduced within the
same paper as his natural deduction. Curien and Herbelin's $\lambda\mu\tilde{\mu}$-calculus \cite{Curien2000:DualityOfComputation}
is one of the languages that, through the lens of the Curry-Howard correspondence,
reveals the computational content of classical $\LK$ proofs.

The sequent calculus' deep symmetries lend themselves to study a variety of
features through duality---such as (co)data-types \cite{Binder2025:Thesis} or
(co)induction \cite{Downen2015:StructuralRecursion}---while at the same time being a suitable intermediate
representation for compilers \cite{Schuster2025:AxCut}.
As such, the study of sequent calculi can guide us in our
exploration of the design of programming languages. Naturally, we may want to apply this
perspective to extensions of traditional logics.

The extensions of logic this thesis is concerned with are derivates of
Girard's \emph{Linear Logic} $\LL$ \cite{Girard19871:LinearLogic}. Term assignments for $\LL$ have also been studied,
both in the intuitionistic \cite{Caires2010:ILLasSessions} and classical setting \cite{Wadler2012:PropositionsAsSessions}.
The languages thus obtained are a form of $\pi$-calculus, with propositions being interpreted as session types
instead of data types. Following the approach explained above, our aim is to
explore the design of concurrent programs by considering a different choice
of sequent calculus as our basis for concurrent computation.
For this purpose, we will first consider the options from which we can choose
when designing the rules of a sequent calculus.

\section{Designing Sequent-Style Calculi}
The meaning of the logical connectives can be understood in terms of the logical rules
which define their justifications or refutations\footnote{This idea was already noted by Gentzen
(\enquote{Die Einführungen stellen sozusagen die \enquote{Definitionen} der betreffenden Zeichen dar,
und die Beseitigungen sind letzten Endes nur Konsequenzen hiervon} \cite[p. 189]{Gentzen1935:Untersuchungen1})
and was further explored by logicians and philosophers such as Wittgenstein \cite{Wittgenstein1974:PhilosophicalGrammar},
Prawitz \cite{Prawitz1974:IdeaOfAGeneralProofTheory}, and Dummett \cite{Dummett1991:LogicalBasisOfMetaphysics}.}.
In a sequent calculus, these rules
can be classified using two criteria: rules either \emph{introduce} their connective
in the conclusion or \emph{eliminate} it from their premises, and they either do so
on the \emph{left} or the \emph{right} side of the turnstile. Thus, for some generic
$n$-ary connective $\odot$, we have four possible combinations as demonstrated
by the following table.

\begin{center}
  \begin{minipage}{0.49\textwidth}
    \centering
    \begin{prooftree}
      \hypo{P_1 \hdots P_n}
      \infer1[]{\Gamma, \odot(T_1, \hdots, T_n) \vdash \Delta}
    \end{prooftree}
    \\ \bigskip
    \begin{prooftree}
      \hypo{\Gamma, \odot(T_1, \hdots, T_n) \vdash \Delta}
      \hypo{P_1 \hdots P_n}
      \infer2[]{C}
    \end{prooftree}
  \end{minipage}\hfill
  \begin{minipage}{0.49\textwidth}
    \centering
    \begin{prooftree}
      \hypo{P_1 \hdots P_n}
      \infer1[]{\Gamma \vdash \odot(T_1, \hdots, T_n), \Delta}
    \end{prooftree}
    \\ \bigskip
    \begin{prooftree}
      \hypo{\Gamma \vdash \odot(T_1, \hdots, T_n), \Delta}
      \hypo{P_1 \hdots P_n}
      \infer2[]{C}
    \end{prooftree}
  \end{minipage}
\end{center}

Here, $P_1, \hdots, P_n$ are premises and $C$ is a conclusion of the form $\Gamma \vdash \Delta$.
The formula $\odot(T_1, \hdots, T_n)$ mentioned in all rules is called the \emph{active formula}.
If it appears in one of the premises, this premise is called the \emph{main premise},
all other premises are called \emph{secondary} or \emph{side premises}.

Concretely, for the conjunction, we may consider the following four rules
which are all valid inferences on formulas of the form $A \land B$.

\begin{center}
  \begin{prooftree}
    \hypo{\Gamma, A, B \vdash \Delta}
    \infer1[\textsc{l-$\land$-i}]{\Gamma, A \land B \vdash \Delta}
  \end{prooftree}
  \qquad
  \begin{prooftree}
    \hypo{\Gamma_1 \vdash A, \Delta_1}
    \hypo{\Gamma_2 \vdash B, \Delta_2}
    \infer2[\textsc{r-$\land$-i}]{\Gamma_1, \Gamma_2 \vdash A \land B, \Delta_1, \Delta_2}
  \end{prooftree}
  \\ \bigskip
  \begin{prooftree}
    \hypo{\Gamma_1, A \land B \vdash \Delta_1}
    \hypo{\Gamma_2 \vdash A, \Delta_2}
    \hypo{\Gamma_3 \vdash B, \Delta_3}
    \infer3[\textsc{l-$\land$-e}]{\Gamma_1, \Gamma_2, \Gamma_3 \vdash \Delta_1, \Delta_2, \Delta_3}
  \end{prooftree}
  \\ \bigskip
  \begin{prooftree}
    \hypo{\Gamma_1 \vdash A \land B, \Delta_1}
    \hypo{\Gamma_2, A, B \vdash \Delta_2}
    \infer2[\textsc{r-$\land$-e}]{\Gamma_1, \Gamma_2 \vdash \Delta_1, \Delta_2}
  \end{prooftree}
\end{center}

However, defining a connective using all four rules is redundant, since some of
these rules can be expressed through each other.
In particular, left-introduction and right-elimination (as well as left-elimination and right-introduction) rules
can be obtained through each other.

This phenomenon is captured by the notion of \emph{bi-expressibility} introduced
by \citet{Ostermann2022:IntroElim}, who demonstrate that the choice of inference rules
gives rise to four functionally complete sub-calculi which correspond to one
of the rows or columns of the table above.

\section{The Computational Content of Proofs}
The choice of sub-calculus does not merely amount to a change in syntax,
it changes our notion of computation.
To demonstrate, we will consider the formula $(A \land B) \to (B \land A)$.

\subsection{Natural Deduction: ND}
Natural deduction only consists of right rules.
The following derivation $\mathcal{D}$ is a proof of the sequent
$\vdash (A \land B) \to (B \land A)$.

\begin{center}
  \begin{prooftree}
    \hypo{A \land B \vdash A \land B}
    \infer1[\textsc{r-$\land$-e$_1$}]{A \land B \vdash B}

    \hypo{A \land B \vdash A \land B}
    \infer1[\textsc{r-$\land$-e$_2$}]{A \land B \vdash A}

    \infer2[\textsc{r-$\land$-i}]{A \land B \vdash B \land A}

    \infer1[\textsc{r-$\to$-i}]{\vdash (A \land B) \to (B \land A)}
  \end{prooftree}
\end{center}

Suppose we had derivations $\mathcal{D}_A$ and $\mathcal{D}_B$ for the sequents
$\vdash A$ and $\vdash B$, respectively, then we can initiate a computation
by applying an implication elimination as follows:

\begin{center}
  \begin{prooftree}
    \hypo{\mathcal{D}}
    \ellipsis{}{\vdash (A \land B) \to (B \land A)}

    \hypo{\mathcal{D}_A}
    \ellipsis{}{\vdash A}
    \hypo{\mathcal{D}_B}
    \ellipsis{}{\vdash B}
    \infer2[\textsc{r-$\land$-i}]{\vdash A \land B}

    \infer2[\textsc{r-$\to$-e}]{\vdash B \land A}
  \end{prooftree}
\end{center}

This proof is not in normal form, in particular, it violates the \emph{subformula property}.

\begin{definition}[Subformula property]
  A derivation $\mathcal{D}$ of a sequent $\Gamma \vdash \Delta$ satisfies
  the subformula property if every formula occurring in $\mathcal{D}$ is
  a subformula of one of the formulas contained in $\Gamma$ or $\Delta$.
\end{definition}

The formula $(A \land B) \to (B \land A)$ occurs in one of its subderivations,
but no longer in the conclusion. The reason for this is an introduction rule
which is immediately followed by an elimination rule. However, we can eliminate
this detour by substituting the proof of $\vdash A \land B$ into $\mathcal{D}$.

\begin{center}
  \begin{prooftree}
    \hypo{\mathcal{D}_A}
    \ellipsis{}{\vdash A}
    \hypo{\mathcal{D}_B}
    \ellipsis{}{\vdash B}
    \infer2[\textsc{r-$\land$-i}]{\vdash A \land B}
    \infer1[\textsc{r-$\land$-e$_1$}]{\vdash B}

    \hypo{\mathcal{D}_A}
    \ellipsis{}{\vdash A}
    \hypo{\mathcal{D}_B}
    \ellipsis{}{\vdash B}
    \infer2[\textsc{r-$\land$-i}]{\vdash A \land B}
    \infer1[\textsc{r-$\land$-e$_2$}]{\vdash A}

    \infer2[\textsc{r-$\land$-i}]{\vdash B \land A}
  \end{prooftree}
\end{center}

This created two new detours: \textsc{r-$\land$-i} followed by \textsc{r-$\land$-e$_1$}
(\textsc{r-$\land$-e$_2$}). Eliminating them finally creates a proof in normal form.

\begin{center}
  \begin{prooftree}
    \hypo{\mathcal{D}_B}
    \ellipsis{}{\vdash B}
    \hypo{\mathcal{D}_A}
    \ellipsis{}{\vdash A}
    \infer2[\textsc{r-$\land$-i}]{\vdash B \land A}
  \end{prooftree}
\end{center}

The proof above corresponds to the (open) $\lambda$-calculus term
\[
  (\lambda x . (proj_2 x, proj_1 x)) (a, b)
\]
and the proof simplifications that were carried out mirror the term's reduction steps
\[
  (\lambda x . (proj_2 x, proj_1 x)) (a, b) \rhd (proj_2 (a, b), proj_1 (a, b)) \rhd (b, a)
\]
At its core, the notion of computation within natural deduction is function creation
and application.

\subsection{Sequent Calculus: LK}
In $\LK$, we no longer have elimination rules on the right, but instead introduction
rules on the left (and right). A derivation $\mathcal{D}$ of the same sequent
now has the form

\begin{center}
  \begin{prooftree}
    \hypo{B \vdash B}
    \hypo{A \vdash A}
    \infer2[\textsc{r-$\land$-i}]{A, B \vdash B \land A}
    \infer1[\textsc{l-$\land$-i}]{A \land B \vdash B \land A}
    \infer1[\textsc{r-$\to$-i}]{\vdash (A \land B) \to (B \land A)}
  \end{prooftree}
\end{center}

To spark a computation, we require the rule of \textsc{cut} (again assuming the existence
of derivations $\mathcal{D}_A$ and $\mathcal{D}_B$).

\begin{center}
  \begin{prooftree}
    \hypo{\mathcal{D}}
    \ellipsis{}{\vdash (A \land B) \to (B \land A)}

    \hypo{\mathcal{D}_A}
    \ellipsis{}{\vdash A}
    \hypo{\mathcal{D}_B}
    \ellipsis{}{\vdash B}
    \infer2[\textsc{r-$\land$-i}]{\vdash A \land B}
    \hypo{B \land A \vdash B \land A}
    \infer2[\textsc{l-$\to$-i}]{(A \land B) \to (B \land A) \vdash B \land A}

    \infer2[\textsc{cut}]{\vdash B \land A}
  \end{prooftree}
\end{center}

A detour in $\LK$ occurs whenever a left and right introduction rule meet in a cut.
We can simplify the cut above by introducing two new ones, both of which have a
principal formula of lower rank.

\begin{center}
  \begin{prooftree}
    \hypo{\mathcal{D}_A}
    \ellipsis{}{\vdash A}
    \hypo{\mathcal{D}_B}
    \ellipsis{}{\vdash B}
    \infer2[\textsc{r-$\land$-i}]{\vdash A \land B}

    \hypo{B \vdash B}
    \hypo{A \vdash A}
    \infer2[\textsc{r-$\land$-i}]{A, B \vdash B \land A}
    \infer1[\textsc{l-$\land$-i}]{A \land B \vdash B \land A}

    \hypo{B \land A \vdash B \land A}
    \infer2[\textsc{cut}]{A \land B \vdash B \land A}

    \infer2[\textsc{cut}]{\vdash B \land A}
  \end{prooftree}
\end{center}

The lower cut can be eliminated via proof composition.
The remaining cut could be eliminated, but a cut with an axiom at the root
does not violate the subformula property.

\begin{center}
  \begin{prooftree}
    \hypo{\mathcal{D}_B}
    \ellipsis{}{\vdash B}
    \hypo{\mathcal{D}_A}
    \ellipsis{}{\vdash A}
    \infer2[\textsc{r-$\land$-i}]{\vdash B \land A}

    \hypo{B \land A \vdash B \land A}
    \infer2[\textsc{cut}]{\vdash B \land A}
  \end{prooftree}
\end{center}

Similarly, we can find a $\lambda\mu\tilde{\mu}$-command, which corresponds
to this proof and its simplifications.
\begin{align*}
  & \; \langle \lambda (x \cdot \alpha). \langle x \,\vert\vert\, Match \{ (x, y) \mapsto \langle (y, x) \,\vert\vert\, \alpha \rangle \} \rangle \,\vert\vert\, (a, b) \cdot \alpha \rangle \\
  \rhd &\; \langle (a, b) \,\vert\vert\, Match \{ (x, y) \mapsto \langle (y, x) \,\vert\vert\, \alpha \rangle \} \rangle \\
  \rhd &\; \langle (b, a) \,\vert\vert\, \alpha \rangle
\end{align*}

Now, opposed to $\ND$, the driver of computation is the interplay between
producers, e.g. an abstraction, and consumers, e.g. a call stack $(a, b) \cdot \alpha$---the
latter not being a first-class entity in $\ND$. Thus, we see that the underlying
logical framework influences not only the structure of our programs, but also
its features and model of computation.

\section{Logical Elimination and Free Deduction}
Neither $\ND$ nor $\LK$ makes use of elimination rules on the left, which is in contrast
to Parigot's \emph{Free Deduction} $\FD$ \cite{Parigot1992:FreeDeduction} where every inference eliminates a formula.
Now, a proof $\mathcal{D}$ of the formula $\tau := (A \land B) \to (B \land A)$ looks as follows.

\begin{center}
  \begin{prooftree}
    \hypo{\tau \vdash \tau}

    \hypo{A \land B \vdash A \land B}
    \hypo{B \land A \vdash B \land A}
    \hypo{B \vdash B}
    \hypo{A \vdash A}
    \infer3[\textsc{l-$\land$-e}]{A, B \vdash B \land A}
    \infer2[\textsc{r-$\land$-e}]{A \land B \vdash B \land A}

    \infer2[\textsc{l-$\to$-e}]{\vdash \tau}
  \end{prooftree}
\end{center}

The derivation already highlights the characteristic of $\FD$'s normal forms: Every
main premise is an axiom which is causing the derivations to grow as right-deep trees.
Thus, as will be explained later, for a computation to occur, an axiom must be subject to both its left and right
elimination rule.

\begin{center}
  \begin{prooftree}
    \hypo{\mathcal{D}}
    \ellipsis{}{\vdash \tau}

    \hypo{A \land B \vdash A \land B}
    \hypo{\mathcal{D}_A}
    \ellipsis{}{\vdash A}
    \hypo{\mathcal{D}_B}
    \ellipsis{}{\vdash B}
    \infer3[\textsc{l-$\land$-e}]{\vdash A \land B}

    \hypo{B \land A \vdash B \land A}

    \infer3[\textsc{r-$\to$-e}]{\vdash B \land A}
  \end{prooftree}
\end{center}

Substituting the proof of $A \land B$ for the axiom $A \land B \vdash A \land B$
yields the first simplification step.

\begin{center}
  \begin{prooftree}[separation=1.0em]
    \hypo{A \land B \vdash A \land B}
    \hypo{\mathcal{D}_A}
    \ellipsis{}{\vdash A}
    \hypo{\mathcal{D}_B}
    \ellipsis{}{\vdash B}
    \infer3[\textsc{l-$\land$-e}]{\vdash A \land B}

    \hypo{B \land A \vdash B \land A}
    \hypo{B \vdash B}
    \infer[no rule]1{A \vdash A}
    \infer2[\textsc{l-$\land$-e}]{A, B \vdash B \land A}
    \infer2[\textsc{r-$\land$-e}]{A \land B \vdash B \land A}
  \end{prooftree}
\end{center}

This reveals a second detour involving \textsc{l-$\land$-e} and \textsc{r-$\land$-e} which
can also be discharged via proof composition to yield the following normal form.

\begin{center}
  \begin{prooftree}
    \hypo{B \land A \vdash B \land A}
    \hypo{\mathcal{D}_B}
    \ellipsis{}{\vdash B}
    \hypo{\mathcal{D}_A}
    \ellipsis{}{\vdash A}
    \infer3[\textsc{l-$\land$-e}]{\vdash B \land A}
  \end{prooftree}
\end{center}

This looks similar to the normal form we reached in $\LK$, but the axiom
$B \land A \vdash B \land A$ cannot be eliminated---and it does not have to be in order to
satisfy the subformula property. Previously, we could interpret
the resulting normal forms as an ordered pair $(b, a)$. This analogy fails in $\FD$
due to the presence of an additional axiom. Therefore, we suggest for $\FD$ to be
interpreted in the context of communicating processes: The axiom $B \land A \vdash B \land A$
creates a binary session channel, and every elimination rule specifies the \enquote{payload}
to send along this channel. Whenever a channel has witnessed communication on both
sides, i.e., a left and right elimination rule meet, we can perform a step of communication.
A normal form having only axioms as main premises thus corresponds to channels
waiting for external communication, i.e., the program only specified the behavior
for one side of the channel. This thesis will provide a suitable term assignment
to make this interpretation of $\FD$ precise.

\section{Contribution and Structure}
We present a term assignment for a version of Parigot's Free Deduction,
called System $\LFD$, and demonstrate how the choice of elimination rules naturally
describes both synchronous and asynchronous communication. Thus showing that $\LFD$ is as an alternative
logical foundation for the $\pi$-calculus. The main technical contribution
of this thesis consists of mechanized proofs for type safety and the
Church-Rosser property for $\LFD$'s term assignment
system, $\piFD$\footnote{The Rocq development is available at \url{https://github.com/flipsim/master-thesis}.
Lemmata and theorems stated in the thesis that have been formalized are annotated
with the symbol $\rocqicon$. For a mapping between the metatheory presented
in this thesis and the formalization, we refer to Appendix \ref{app:mechanization-map}.}.

The rest of this thesis is structured in three parts.
\begin{itemize}
  \item The remainder of the first part gives an introduction into the logical system
        of Free Deduction and how $\LFD$ is derived from it.
  \item The second part introduces the term assignment system $\piFD$ and its relation
        to existing Curry-Howard interpretations of linear logics.
  \item The last part delves into the metatheory of $\piFD$: We prove type
        preservation (Theorem \ref{theorem:preservation}), deadlock freedom
        (Theorem \ref{theorem:progress}), and the Church-Rosser property
        (Theorem \ref{theorem:church-rosser}). To formalize Theorem \ref{theorem:progress},
        we identify the set of terminated processes---which coincides with (weak) normal forms
        of the underlying logical framework---and show that they are irreducible (Lemma \ref{lemma:normalforms}).
\end{itemize}